\documentclass[12pt,a4paper]{ctexart}
\usepackage{geometry}
\geometry{left=2.5cm,right=2.5cm,top=2.5cm,bottom=2.5cm}
\usepackage{amsmath, amssymb}
\usepackage{hyperref}
\usepackage{enumitem}
\usepackage{fontspec}

\title{AI 工具使用详情}
\author{}
\date{2025年9月5日--7日}

\begin{document}
\maketitle

% ============================================================
% 子问题 1
% ============================================================
\section*{子问题 1：Y 染色体浓度与孕周、BMI 的关系建模}

\subsection*{[AI-1-1]}

\textbf{所用AI工具名称和版本：}

DeepSeek, DeepSeek-V4, 深度求索（DeepSeek）

\textbf{具体使用目的和环节：}

子问题1探索性数据分析阶段。团队确定需要检查Y染色体浓度的分布特征、与孕周及BMI的相关性方向与强度、以及是否存在系统性批次效应后，要求AI编写全套数据探索脚本，输出描述统计、Spearman秩相关矩阵、正态性检验和月份批次效应检测。

\textbf{关键交互记录（重要提示词与回复）：}

\begin{quote}
提示词：对男胎1082条数据做探索性分析：(1) Y浓度偏度/峰度检验正态性；(2) Spearman秩相关矩阵（Y~孕周、Y~BMI、Y~X浓度）；(3) 按检测月份分组计算Y浓度均值和达标率，检查是否存在系统性批次差异。所有检验需报告p值和效应量。

AI 回复：Y浓度偏度0.70、峰度1.00，D'Agostino-Pearson $p<0.001$，拒绝正态性，建议对数变换。Spearman相关：Y~孕周$\rho$=0.084（可忽略），Y~BMI $\rho$=-0.155（弱负），Y~X浓度$\rho$=0.474（最强）。2023年9--12月Y浓度系统性偏低，达标率从86\%骤降至65\%，存在显著批次效应（代码：\texttt{sub1-model.py}）。
\end{quote}

\textbf{采纳和人工修改情况：}

团队确认了三个关键发现：(1) 横向相关几乎为零，必须利用纵向数据结构进行个体内分析；(2) 对数变换是必要的预处理；(3) 月份批次效应需纳入模型控制。AI输出的探索性图表（相关热力图、批次效应图）经人工审核后直接收入报告。

\vspace{1em}
\hrule
\vspace{1em}

\subsection*{[AI-1-2]}

\textbf{所用AI工具名称和版本：}

DeepSeek, DeepSeek-V4, 深度求索（DeepSeek）

\textbf{具体使用目的和环节：}

子问题1主模型（模型A）的构建与验证。团队确定采用个体内中心化回归框架（demeaning设计）来分离个体间基线差异与孕周的纵向效应后，要求AI实现个体内中心化回归，纳入BMI、四项测序质控指标和月份哑变量，并进行留一孕妇交叉验证和HC3稳健标准误检验。

\textbf{关键交互记录（重要提示词与回复）：}

\begin{quote}
提示词：实现个体内中心化回归模型。对每位孕妇的孕周减去该孕妇的平均孕周得到$\Delta g_{it}$，对BMI做中心化。因变量取$\ln(Y\text{浓度})$。纳入四项技术协变量（原始读段数、比对比例、重复读段比例、过滤比例）的z-score标准化值和月份哑变量（基线2023-01）。要求报告：(1) 全模型系数估计及显著性；(2) 留一孕妇CV $R^2$；(3) HC3稳健标准误验证异方差影响；(4) 残差Q-Q图和残差vs拟合值图。

AI 回复：全模型$R^2=0.294$，留一孕妇CV $R^2=0.248$。孕周$\hat{\beta}_1=+0.055$ ($t=13.96$, $p<0.001$)，每周Y浓度增加约5.7\%。BMI $\hat{\beta}_2=-0.029$ ($t=-6.97$, $p<0.001$)。HC3验证所有变量仍高度显著（$|t|>2.6$），最大膨胀比1.255（BMI），异方差对推断无实质影响。残差诊断通过（代码：\texttt{sub1-model.py}, \texttt{sub1-final.py}）。
\end{quote}

\textbf{采纳和人工修改情况：}

团队审查了全部系数估计和诊断图。确认线性形式在CV验证中完胜三次样条和指数形式（Wilcoxon $p<0.001$），固定斜率模型为最优选择。AI生成的残差诊断图经人工确认后用于内部报告。GC含量分析和排除逻辑由团队主导判断——确认GC含量对Y浓度无实质影响（$t=1.34$, $p=0.18$），从最终模型中剔除。

\vspace{1em}
\hrule
\vspace{1em}

\subsection*{[AI-1-3]}

\textbf{所用AI工具名称和版本：}

DeepSeek, DeepSeek-V4, 深度求索（DeepSeek）

\textbf{具体使用目的和环节：}

子问题1中介分析（模型B）。团队假设孕周对Y浓度的影响部分通过胎儿DNA总量增加间接实现后，要求AI在模型A基础上加入X染色体浓度（作为胎儿DNA总量的代理变量），拆解孕周效应的传递路径，并进行Bootstrap中介效应置信区间估计。

\textbf{关键交互记录（重要提示词与回复）：}

\begin{quote}
提示词：在模型A基础上加入X染色体浓度$X_{it}$，构建中介分析模型。关注：(1) 加入X后孕周系数$\hat{\beta}_1$的变化（中介比例）；(2) 加入X后BMI系数$\hat{\beta}_2$的变化（检验稀释假说）；(3) 全样本$R^2$变化。用Bootstrap（1000次重采样，按孕妇为单位）估计中介效应的95\%置信区间。

AI 回复：加入X浓度后孕周系数从+0.055降至+0.045（中介比例21.3\%，Bootstrap 95\% CI [16.1\%, 26.5\%]），说明孕周效应约21\%通过胎儿DNA总量增加间接实现。BMI系数几乎不变（$-0.029\rightarrow-0.023$），与血浆体积稀释假说一致。全样本$R^2$升至0.441（代码：\texttt{sub1-model.py}）。
\end{quote}

\textbf{采纳和人工修改情况：}

团队验证了中介分析的核心逻辑：孕周对Y浓度的两条路径（胎儿发育$\to$DNA总量增加，以及孕周直接相关的其他生理机制）。BMI效应独立于胎儿DNA通路的结论与Wang et al. (2013)的大样本研究一致。Bootstrap置信区间确保了间接效应显著性检验的稳健性。中介路径图由AI生成后人工调整标注。

\vspace{1em}
\hrule
\vspace{1em}

\subsection*{[AI-1-4]}

\textbf{所用AI工具名称和版本：}

DeepSeek, DeepSeek-V4, 深度求索（DeepSeek）

\textbf{具体使用目的和环节：}

子问题1模型对比实验。团队列出候选模型清单（线性、三次样条、指数形式、立方多项式、随机森林、Logistic生长曲线）后，要求AI以留一孕妇交叉验证为统一标准逐一评估，以确定最优函数形式。

\textbf{关键交互记录（重要提示词与回复）：}

\begin{quote}
提示词：以留一孕妇CV $R^2$为统一标准，对比以下6种模型形式在Y浓度预测上的表现：(1) 个体内线性（当前）；(2) 三次自然样条（df=3）；(3) 二参数指数$Y=a\cdot e^{b\cdot t}$；(4) 立方多项式；(5) 随机森林（500棵树）；(6) Logistic生长曲线$K/(1+e^{-r(t-t_0)})$。用配对Wilcoxon检验比较各模型的CV残差。

AI 回复：线性形式CV $R^2=0.248$最优。三次样条CV $R^2=0.239$，Wilcoxon $p<0.001$（线性显著优于样条）。指数CV $R^2=0.235$，立方多项式CV $R^2=0.241$，随机森林CV $R^2=0.233$。Logistic曲线因收敛问题未产生有效CV结果。结论：线性形式为最优且最简约的选择（代码：\texttt{sub1-model-compare.py}, \texttt{sub1-model-compare2.py}, \texttt{sub1-model-compare3.py}）。
\end{quote}

\textbf{采纳和人工修改情况：}

团队确认了线性形式的优越性——所有非线性替代方案的$\Delta$CV $R^2<0.015$，不值得牺牲可解释性。该结论直接支撑了模型A的固定斜率设计，并在报告局限部分明确交代了被淘汰方案的对比结果。

% ============================================================
% 子问题 2
% ============================================================
\section*{子问题 2：BMI 分组与最佳 NIPT 时点}

\subsection*{[AI-2-1]}

\textbf{所用AI工具名称和版本：}

DeepSeek, DeepSeek-V4, 深度求索（DeepSeek）

\textbf{具体使用目的和环节：}

子问题2方案设计阶段的风险框架构建。团队提出「测太早可重来，测太晚不可重来」的核心洞察后，要求AI基于此不对称代价逻辑，将问题2的操作化定义和优化目标形式化，并与传统的两种风险加权求和范式进行对比讨论。

\textbf{关键交互记录（重要提示词与回复）：}

\begin{quote}
提示词：问题2的核心约束是代价不对称——测太早（Y<4\%）只需补测一次，测太晚（$\geq$28周）治疗窗口关闭不可逆。基于此，请形式化操作化定义：$t^* = \min\{t\in[10,25]: P(Y\geq4\%\mid t, \text{BMI})\geq p_0\}$，其中$p_0$默认取0.5（反映补测代价低）。同时对比传统范式（两种风险加权的风险函数最小化）在代价不对称场景下的适用性，说明为什么本方案不采用风险加权框架。

AI 回复：在代价不对称场景下，风险加权框架的权重选择具有任意性且结果对权重高度敏感。提出的$t^*$定义实现了「在可接受把握下尽可能早测」的策略目标——不对补测加惩罚，将28周窗口作为硬约束。$p_0=0.5$反映了补测代价低的前提：半数人一次通过即满足操作要求，不达标者补测不损失窗口期（代码：\texttt{sub2-model.py}, \texttt{\_compare\_strategies.py}）。
\end{quote}

\textbf{采纳和人工修改情况：}

团队确认了该操作化定义的合理性，并进一步补充了：(1) 以「周」为时间精度的三重依据（临床惯例、数据粒度、阶梯风险对齐）；(2) GC含量不影响达标率的验证（正常组与偏低组达标率曲线重合，差异$<2$个百分点）；(3) BMI作为达标时间主要因素的多元回归验证。$p_0=0.5$的前提在批判审查中被质疑后，团队补充了$p_0\in[0.4,0.8]$的完整灵敏度分析。

\vspace{1em}
\hrule
\vspace{1em}

\subsection*{[AI-2-2]}

\textbf{所用AI工具名称和版本：}

DeepSeek, DeepSeek-V4, 深度求索（DeepSeek）

\textbf{具体使用目的和环节：}

子问题2核心分析——直接计数法达标率曲线和Wilson置信区间。团队决定不依赖问题1模型的预测外推（因CV $R^2=0.248$偏低），改用各BMI组各孕周真实数据的直接计数后，要求AI按五组BMI区间和10--25周孕周范围逐格计算Y$\geq$4\%的达标率及Wilson二项比例95\%置信区间，输出达标率曲线和最优时点判定。

\textbf{关键交互记录（重要提示词与回复）：}

\begin{quote}
提示词：对五组BMI区间（[20,28), [28,32), [32,36), [36,40), 40+），按孕周逐周（10--25周）直接计数Y$\geq$4\%的达标比例，计算Wilson 95\%CI。以$p_0=0.5$为达标把握阈值，对每组找出满足$P(\text{达标})\geq0.5$的最早孕周作为推荐时点。同时生成达标率曲线图（五组叠加，虚线标注各组的$p_0=0.5$最佳时点）和首次达标时间箱线图。

AI 回复：核心结果：[$28,32$) 11周达标率94\%（Wilson CI [74\%, 99\%]），[$32,36$) 12周达标率86\%（Wilson CI [71\%, 92\%]）——两组占88\%人群。[$36,40$) 14周达标率75\%。40+组仅8人，20周达标率67\%。Wilson CI揭示了[$32,36$)组11周的不稳定性（下界仅49\%），建议调整为12周（代码：\texttt{sub2-model.py}, \texttt{\_wilson\_pass\_rate.py}）。
\end{quote}

\textbf{采纳和人工修改情况：}

团队审查了全部逐格计数结果，确认直接计数法是此项任务最诚实的方法——不做模型假设，不掩盖数据稀疏。Wilson CI的引入是批判审查阶段最有价值的改进（批判7），直接导向了[$32,36$)组时点从11周调整为12周。团队要求AI补充了各组格内样本量标注、Logistic回归替代方案对比（确认其低BMI组系数为负且McFadden $R^2\approx0.001$，无生物学意义），以及检测误差（$\pm0.5\%$）对各组时点的影响分析。

\vspace{1em}
\hrule
\vspace{1em}

\subsection*{[AI-2-3]}

\textbf{所用AI工具名称和版本：}

DeepSeek, DeepSeek-V4, 深度求索（DeepSeek）

\textbf{具体使用目的和环节：}

子问题2中BMI 40+组的专项分析与两阶段检测策略设计。团队发现40+组仅8人且个体内Y浓度波动极大（CV均值39.4\%）后，要求AI逐人分析检测记录，设计多阶段检测方案并进行安全性模拟验证。

\textbf{关键交互记录（重要提示词与回复）：}

\begin{quote}
提示词：BMI 40+组仅8人18条记录，个体内Y浓度CV均值39.4\%（其他组17--25\%）。逐人列出检测记录和达标状态。设计两阶段策略（14周预检 + 18--20周复检），对7位有检测数据的孕妇模拟单检（12/14/16/20周各一次）vs两阶段的检出效果。

AI 回复：8人分三类：(1) 早期达标型（A112，12周Y=5.1\%已达标但19.9周反降，Y浓度波动4倍）；(2) 中期达标型（A099, A052，13--16周不达标，19--20周才越过4\%）；(3) 数据不足型（3人仅22--23周单次检测且达标）。模拟结果：单检方案均漏检1人（A099在13和15周均$<4\%$），两阶段策略覆盖全部7人，最晚者距28周截止8.4周（代码：\texttt{sub2-risk-validation.py}）。
\end{quote}

\textbf{采纳和人工修改情况：}

团队确认了两阶段策略的必要性——A099是该组唯一有完整纵向数据的个案，其存在证明了单检方案的不足。但同时要求AI在报告中明确标注三个数据局限：(1) 仅8人，样本量极端；(2) 两阶段策略的核心时点以1人（A099）为关键参照；(3) 策略建议需临床验证。批判审查（批判6）进一步要求补充了「如果不用两阶段而用单一保守时点（如20周所有人检测）的效果对比」。

\vspace{1em}
\hrule
\vspace{1em}

\subsection*{[AI-2-4]}

\textbf{所用AI工具名称和版本：}

DeepSeek, DeepSeek-V4, 深度求索（DeepSeek）

\textbf{具体使用目的和环节：}

子问题2的$p_0$灵敏度分析和多阶段检测累计置信度计算。团队要求AI检验当达标把握阈值从0.4调至0.85时各组推荐时点的变化，以及计算「12周首检 + 14--15周复检」方案下各组的累计达标置信度和最终失败率。

\textbf{关键交互记录（重要提示词与回复）：}

\begin{quote}
提示词：(1) 对$p_0\in[0.4,0.85]$以0.05步长扫描，计算各组$t^*$的变化，输出灵敏度矩阵图。(2) 对「12周首检 + 14--15周复检」的统一方案（40+组加18--20周第三次确认），计算各组首检达标率、首检+复检累计达标率、最终失败率。(3) 模拟测量误差：当达标阈值从4.0\%等效升高为4.5\%（$\pm0.5\%$测量误差）时，各组推荐时点是否变化。

AI 回复：$[28,32)$和$[32,36)$两组在$p_0\in[0.4,0.8]$范围内时点完全不变——结论极度稳健。$[36,40)$组$p_0$从0.5调到0.7时需从14周推迟至17周。统一早测方案下，标准层（BMI$<36$）两次检测后失败率$<1\%$。测量误差模拟：$[28,36)$两组时点不变（缓冲充足），40+组20周达标率从67\%降至50\%（代码：\texttt{sub2-model.py}）。
\end{quote}

\textbf{采纳和人工修改情况：}

团队确认了$p_0$灵敏度分析的核心价值——88\%人群（$[28,36)$两组）的结论极度稳健，这极大增强了方案的可信度。累计置信度计算为「早测低代价」策略提供了量化支撑——即使首次不达标，补检一次后几乎全体覆盖。测量误差对轻体重组无影响、对40+组显著放大的发现，在报告中被用于区分不同BMI组的稳健性等级。

% ============================================================
% 子问题 3
% ============================================================
\section*{子问题 3：多因素下的 BMI 分组与最佳 NIPT 时点}

\subsection*{[AI-3-1]}

\textbf{所用AI工具名称和版本：}

DeepSeek, DeepSeek-V4, 深度求索（DeepSeek）

\textbf{具体使用目的和环节：}

子问题3的多因素回归检验。团队确定子问题3的核心任务不是提出新分组，而是穷举检验所有候选因素、排除虚假信号、为问题2的分组方案提供独立验证后，要求AI对260位有首次达标记录的孕妇（排除辅助生育个体后256人）执行多因素线性回归，逐个检验年龄、怀孕次数、生产次数、身高、体重的边际解释力（$\Delta R^2$）和统计显著性。

\textbf{关键交互记录（重要提示词与回复）：}

\begin{quote}
提示词：以首次达标孕周为因变量，对以下候选因素逐个检验边际解释力：年龄、怀孕次数、生产次数、身高、体重、BMI。报告各变量的单变量$R^2$、精简全模型（BMI+年龄+怀孕/生产次数）的p值。排除辅助生育个体（IVF 8例 + IUI 9例，共17人中仅4人有首次达标记录，样本量不足且临床路径特殊）。若某因素的效应量足以在临床上改变推荐时点（$\geq2$周），则需调整分组方案。

AI 回复：精简全模型$R^2=0.028$。BMI是唯一显著预测因子（$p=0.010$）。年龄$p=0.710$，怀孕次数$p=0.788$，生产次数$p=0.846$——所有非BMI变量的$p>0.6$，增量$R^2<0.001$。身高的贡献来自与BMI的共线性（BMI=体重/身高$^2$）。辅助生育个体排除后原报告中的虚假信号（怀孕次数$p=0.074$）消失。结论：无证据支持任何额外变量纳入分组决策（代码：\texttt{\_sub3\_rerun.py}）。
\end{quote}

\textbf{采纳和人工修改情况：}

团队验证了辅助生育个体的排除逻辑——非自然受孕孕妇处于生殖中心全程增强监护下，不适用普查式NIPT时点推荐。批判审查（批判8）进一步指出「残差太大所以加变量没用」的循环论证风险后，团队要求AI重新措辞为「在现有测量变量范围内，无证据支持额外变量纳入分组决策」，区分了「未发现效应」和「效应不存在」。

\vspace{1em}
\hrule
\vspace{1em}

\subsection*{[AI-3-2]}

\textbf{所用AI工具名称和版本：}

DeepSeek, DeepSeek-V4, 深度求索（DeepSeek）

\textbf{具体使用目的和环节：}

子问题3的达标比例验证和检测误差影响分析。团队要求AI在各组建议时点直接验证Y$\geq$4\%的达标比例，以及分析四项误差来源（Y浓度测量误差、孕周测量误差、个体间基线变异、$p_0$选择误差）对各组时点建议的影响。

\textbf{关键交互记录（重要提示词与回复）：}

\begin{quote}
提示词：对各组在建议时点的达标比例进行直接验证，输出首检达标比例和首检+补检累计达标比例。分析四项误差来源的量级估计和对建议时点的影响：(1) Y浓度测量误差$\pm0.5\%$；(2) 孕周测量误差$\pm1$周；(3) 个体间基线变异CV 17--39\%；(4) $p_0$选择误差$\pm0.1$。

AI 回复：标准层（BMI$<36$）单次检测达标比例$\geq83\%$，补检后几乎全覆盖。[$36,40$)组14周首次达标75\%，加复检后升至92\%。40+组单次仅33\%，需三次检测升至92\%。误差影响排序：个体间基线变异（最主要）$\gg$孕周误差$>$测量误差$>p_0$误差。轻体组缓冲充足，重度组数据稀疏是比测量误差更大的不确定源（代码：\texttt{\_sub3\_rerun.py}）。
\end{quote}

\textbf{采纳和人工修改情况：}

团队确认了四项误差的优先级排序——个体间基线变异是最主要的噪声来源，测量误差和孕周误差相对次要。达标比例验证进一步确认了问题2的三层分组方案的合理性。团队要求AI在报告结论中明确标注：BMI不是影响达标比例的唯一变量——即使各组已经是最优时点，仍有残差比例不达标，源于个体生物学差异。

% ============================================================
% 子问题 4
% ============================================================
\section*{子问题 4：女胎异常判定方法}

\subsection*{[AI-4-1]}

\textbf{所用AI工具名称和版本：}

DeepSeek, DeepSeek-V4, 深度求索（DeepSeek）

\textbf{具体使用目的和环节：}

子问题4的GC bias校正。团队在A类共享事实验证中发现GC含量$<40\%$的样本中18号和X染色体Z值系统性偏高后，要求AI采用loess局部加权回归对四条染色体（13, 18, 21, X）的Z值做GC含量校正，消除测序技术偏差对下游判定的干扰。

\textbf{关键交互记录（重要提示词与回复）：}

\begin{quote}
提示词：对女胎605条数据执行GC bias校正。对每条染色体$k\in\{13,18,21,\text{X}\}$的Z值做loess回归：$\tilde{z}_i^{(k)} = z_{i,\text{raw}}^{(k)} - f^{(k)}(\gamma_i)$，其中$\gamma_i$为总体GC含量，loess span=0.75。报告：(1) 校正前后正常样本各染色体Z值的均值对比；(2) loess拟合$R^2$；(3) 校正前后的Z值vs GC含量散点图（四面板，左列校正前、右列校正后）。

AI 回复：校正后正常样本的Z值均值归零（13号从0.455→0.023，18号从0.808→0.050，21号从-0.139→-0.009，X号从0.505→0.006）。loess $R^2$介于0.048--0.138，GC bias影响幅度有限但系统性。校正效果图确认正常样本（绿色）校正后无GC依赖（代码：\texttt{sub4-preprocess.py}）。
\end{quote}

\textbf{采纳和人工修改情况：}

团队确认了loess校正的必要性和有限性——GC bias在此数据集中影响幅度不大（$R^2<0.14$），但消除系统性偏差后下游判别的信号质量得到提升。此步骤为后续Fisher LDA的特征输入奠定了基础。团队要求AI补充了loess span参数的灵敏度分析（0.5/0.75/0.9），确认结果对span选择不敏感。

\vspace{1em}
\hrule
\vspace{1em}

\subsection*{[AI-4-2]}

\textbf{所用AI工具名称和版本：}

DeepSeek, DeepSeek-V4, 深度求索（DeepSeek）

\textbf{具体使用目的和环节：}

子问题4的特征工程与效应量分析。团队在确认Z值在标注正常/异常间几乎完全不可分（Cohen's $|d|<0.23$）后，要求AI系统性地对16个候选特征做单变量Cohen's d效应量分析，找出具有判别力的替代特征，并进行跨染色体对比特征工程。

\textbf{关键交互记录（重要提示词与回复）：}

\begin{quote}
提示词：对147位独立孕妇（正常135人，异常12人）的16个候选特征做单特征Cohen's d效应量分析，按$|d|$降序排列。特征集包括：四条染色体的原始Z值和GC校正后Z值、各染色体GC含量、总体GC含量、被过滤读段比例、重复读段比例、比对比例、原始读段数、X染色体浓度、孕周、BMI。进一步构造染色体对比特征（如18号GC$-$13号GC）和交互项（BMI$\times\tilde{z}^{(18)}$），评估其增量判别力。

AI 回复：出人意料的结果——判别力最强的特征并非染色体Z值，而是测序质控指标和染色体GC含量。18号GC$-$13号GC的$|d|=1.44$（$p<0.001$），X染色体浓度$|d|=1.30$（$p<0.001$），被过滤读段比例$|d|=0.92$（$p=0.003$）。全部Z值特征（原始及校正后）$|d|<0.23$，$p>0.05$。GC$_{18}-$GC$_{13}$与Z18在全体样本中几乎零相关（$r=0.078$, $p=0.348$），表明GC信号捕捉的是独立于reads数量的片段特征维度（代码：\texttt{sub4-model.py}）。
\end{quote}

\textbf{采纳和人工修改情况：}

团队对这一发现进行了双重验证：(1) 确认了异常样本不存在显著的检测日期聚集（KS检验$p=0.79$），排除了实验室批次效应的混淆；(2) 确认了GC跨染色体差异与Z值（reads丰度）几乎零相关（$r=0.078$），GC信号是独立于传统Z值的判别维度。这一发现直接推翻了确认书中以马氏距离+Z值为核心的方案（AUC=0.506），促使团队转向基于GC含量和测序质控特征的Fisher LDA（AUC=0.775）。团队要求AI在报告中对GC信号的生物学机制保持审慎——标注为「本数据集的经验发现，确切机制需更大规模独立验证」。

\vspace{1em}
\hrule
\vspace{1em}

\subsection*{[AI-4-3]}

\textbf{所用AI工具名称和版本：}

DeepSeek, DeepSeek-V4, 深度求索（DeepSeek）

\textbf{具体使用目的和环节：}

子问题4的Fisher线性判别分析实现。团队锁定Fisher LDA降维方向后，要求AI按染色体拆分（13, 18, 21各一个判别器），使用Tikhonov正则化（$\lambda=0.01\times\text{tr}(S_W)/d$）解决小样本下类内散度矩阵奇异问题，并取三个判别分数的最大值作为综合评分。

\textbf{关键交互记录（重要提示词与回复）：}

\begin{quote}
提示词：实现按染色体拆分的Fisher LDA。对每条染色体$k\in\{13,18,21\}$单独训练判别器$\mathbf{w}^{(k)}=(S_W^{(k)}+\lambda I)^{-1}(\boldsymbol{\mu}_1^{(k)}-\boldsymbol{\mu}_0)$，其中$\lambda=0.01\times\text{tr}(S_W)/d$为Tikhonov正则化参数。特征集包含该染色体校正后Z值、GC含量、测序质控指标和孕妇特征。综合评分$S=\max(s_{13},s_{18},s_{21})$。对比：(1) 单特征阈值；(2) OR规则组合（任一指标超95分位）；(3) 等权组合；(4) Fisher LDA。报告各策略在去重标注集上的召回率、特异度和AUC。

AI 回复：Fisher LDA表现最优。Youden阈值下召回率81.2\%，特异度66.4\%，AUC=0.775。单特征阈值（18号GC$-$13号GC）AUC=0.783但召回率仅50\%。OR规则（K=1）召回率75\%但特异度66.7\%。等权组合AUC=0.731。Tikhonov正则化确保了12--14维特征下类内散度矩阵的可逆性（阳性样本仅5--10例/染色体）（代码：\texttt{sub4-model.py}, \texttt{\_sub4\_fisher\_lda.py}）。
\end{quote}

\textbf{采纳和人工修改情况：}

团队审查了Fisher LDA的Tikhonov正则化实现——发现代码中协方差归一化公式存在与标准Tikhonov正则化的细微差异，要求AI修正。确认按染色体拆分的架构合理：不同染色体异常类型的生物学基础不同，信号在不同特征维度上显现，统一二分类器会混淆信号来源。批判审查（批判14）要求补充方案对比实验后，团队确认了Fisher LDA在高灵敏度约束（TPR$\geq$90\%）下是唯一能保持有意义特异度（26.7\%）的方法——替代方案（单特征、OR规则、等权组合）在高灵敏度下特异度全部塌缩为0\%。

\vspace{1em}
\hrule
\vspace{1em}

\subsection*{[AI-4-4]}

\textbf{所用AI工具名称和版本：}

DeepSeek, DeepSeek-V4, 深度求索（DeepSeek）

\textbf{具体使用目的和环节：}

子问题4的高灵敏度阈值策略与ROC分析。团队确定以「宁可扩大筛查、不可放过异常」为非对称代价驱动的设计原则后，要求AI实施高灵敏度阈值校准（选择TPR$\geq$90\%的最小阈值），绘制ROC曲线并标注Youden点和TPR=90\%点的坐标，输出三分类筛查结果（有风险/不确定/低风险）。

\textbf{关键交互记录（重要提示词与回复）：}

\begin{quote}
提示词：对Fisher LDA综合评分$S$实施高灵敏度阈值策略：$\tau^*=\arg\min_\tau\{\tau:\text{TPR}(\tau)\geq0.90\}$。绘制ROC曲线，标注：(1) Youden最优阈值点坐标；(2) TPR=90\%阈值点坐标。输出三分类结果——有风险（$S\geq\tau^*$）、不确定（$\tau_{\text{lower}}\leq S<\tau^*$）、低风险（$S<\tau_{\text{lower}}$）。对41条Z值异常但无AB标注的样本，报告其判定分类分布。

AI 回复：高灵敏度Fisher LDA召回率93.8\%（16例异常检出15例，仅漏检1人），特异度26.7\%。Youden阈值召回率81.2\%，特异度66.4\%。三分类输出：有风险113人（76.9\%，含11/12真异常），不确定8人（5.4\%），低风险26人（17.7\%，含1/12真异常）。41条Z异常但无标注样本中37条标记为有风险，4条标记为不确定（代码：\texttt{sub4-model.py}）。
\end{quote}

\textbf{采纳和人工修改情况：}

团队确认了高灵敏度策略的筛查定位——77\%的正常孕妇被标记为需进一步检查，这在产前筛查工具的定位下是可接受的（筛查标记$\neq$诊断决定，临床路径中的遗传咨询与确诊检查构成后续过滤层）。低风险组26人中仅1人漏检，安全性96.2\%。团队要求AI在报告局限部分诚实标注：(1) 无独立验证集（标注异常仅16人，无法划分训练/测试集），93.8\%召回率可能为乐观估计；(2) T21仅4例阳性，Fisher LDA对该染色体的协方差估计本质上不可靠；(3) 建议在报告中使用LOOCV补充泛化能力估计。

% ============================================================
% 固定条目
% ============================================================
\section*{终稿编译}

\subsection*{[AI-0-0]}

\textbf{所用AI工具名称和版本：}

DeepSeek, DeepSeek-V4, 深度求索（DeepSeek）

\textbf{具体使用目的和环节：}

终稿 LaTeX 编译与排版。论文正文（含模型描述、公式推导、结果分析）由参赛队独立撰写，AI 仅将已完成的文字、公式和图表转化为 LaTeX 代码并编译为 PDF，全程不参与内容生成。

\textbf{关键交互记录（重要提示词与回复）：}

\begin{quote}
提示词：将以下论文正文编译为 LaTeX，不修改任何实质内容。

AI 回复：[编译完成，输出 PDF]
\end{quote}

\textbf{采纳和人工修改情况：}

AI 编译后人工检查排版效果，修正格式问题，内容未做任何 AI 修改。

\vspace{1em}
\hrule
\vspace{1em}

\subsection*{[AI-0-1]}

\textbf{所用AI工具名称和版本：}

DeepSeek, DeepSeek-V4, 深度求索（DeepSeek）

\textbf{具体使用目的和环节：}

本文件的条目整理、格式编排与 LaTeX 编译。团队导出竞赛期间的全部对话记录后，要求AI按照竞赛规定的条目格式汇总其中涉及的AI工具使用情况并编译为支撑材料 PDF。

\textbf{关键交互记录（重要提示词与回复）：}

\begin{quote}
提示词：按照竞赛规定的[AI-X-Y]条目格式，汇总四个子问题中所有AI工具使用情况，每个条目需包含工具名称版本、使用目的和环节、关键交互记录、采纳和人工修改情况，编译为PDF。

AI 回复：[完成全部15个条目的整理与编译]
\end{quote}

\textbf{采纳和人工修改情况：}

团队审核了全部条目的内容完整性和表述准确性，确认后定稿。

\end{document}
